\exercisesection{Resolutions}

\begin{enumerate}
\def\labelenumi{\arabic{enumi})}
\tightlist
\item
  Let \(a, b\) two elements of A, such that the left annihilator of \(a\) (resp. \(b\) ) is the ideal \(Ab\) (resp. \(Aa\) ). Show that the sequence
\end{enumerate}

\[
\dots \longrightarrow \mathrm{A} _ {s} \xrightarrow {\delta_ {a}} \mathrm{A} _ {s} \xrightarrow {\delta_ {b}} \mathrm{A} _ {s} \xrightarrow {\delta_ {a}} \mathrm{A} _ {s} \longrightarrow \dots \longrightarrow \mathrm{A} _ {s} \xrightarrow {\delta_ {a}} \mathrm{A} _ {s} \longrightarrow \mathrm{A} / \mathrm{A} a \longrightarrow 0,
\]

where \(\delta_{a}\) (resp. \(\delta_{b}\) ) denotes right multiplication by \(a\) (resp. \(b\) ), defines a free resolution of the A-module A/Aa. Show that this construction applies in the following two cases:

\(\alpha)\) Let \(\mathbf{A}_0\) a ring; we take \(\mathrm{A} = \mathrm{A}_0(\varepsilon)\) (X, p. 27), \(a = b = \varepsilon\) .

β) Let G be a finite cyclic group, σ a generator of G, k a commutative ring. We take

\[
\mathbf {A} = k ^ {(\mathbf {G})}, \quad a = 1 - e _ {\sigma}, \quad b = \sum_ {g \in \mathbf {G}} e _ {g}.
\]

\begin{enumerate}
\def\labelenumi{\arabic{enumi})}
\setcounter{enumi}{1}
\tightlist
\item
  We assume the ring A is commutative. Let B be an A-algebra; for \(n \geqslant 0\) , we denote by \(\mathcal{A}^{n}(\mathbf{B})\) the tensor product over A of \((n + 1)\) modules equal to B. We set \(\mathcal{A}^{n}(\mathbf{B}) = 0\) for \(n < 0\) and \(\mathcal{A}(\mathbf{B}) = \bigoplus \mathcal{A}^{n}(\mathbf{B})\) We define a graded A-endomorphism of ascending degree +1 of \(\mathcal{A}(\bar{\mathbf{B}})\) by setting:
\end{enumerate}

\[
d (b _ {0} \otimes \dots \otimes b _ {n}) = \sum_ {i = 0} ^ {n + 1} (- 1) ^ {i} b _ {0} \otimes \dots \otimes b _ {i - 1} \otimes 1 _ {\mathbf {B}} \otimes b _ {i} \otimes \dots \otimes b _ {n} \quad \text { for } b _ {0}, \dots , b _ {n} \in \mathbf {B}.
\]

\begin{enumerate}
\def\labelenumi{\alph{enumi})}
\tightlist
\item
  Show that \(d \circ d = 0\) , so that \((\mathcal{A}(\mathbf{B}), d)\) is a complex of A-modules. If M is an A-module, we denote by \(\mathcal{A}(\mathbf{B}, \mathbf{M})\) the complex such that \(\mathcal{A}^{n}(\mathbf{B}, \mathbf{M}) = \mathcal{A}^{n}(\mathbf{B}) \otimes_{\mathbf{A}} \mathbf{M}\) , with differential \(d \otimes 1_{\mathbf{M}}\) . b) Suppose that the algebra B is a faithfully flat A-module (X, p. 169, exercise 9). Show that the complex \(\mathcal{A}(\mathbf{B}, \mathbf{M})\) defines a right resolution of the A-module M. (It suffices to verify that the complex \(\mathbf{B} \otimes_{\mathbf{A}} \mathcal{A}(\mathbf{B}, \mathbf{M})\) defines a resolution of \(\mathbf{B} \otimes_{\mathbf{A}} \mathbf{M}\) ; construct an A-linear homotopism between these two complexes.)
\end{enumerate}

\begin{enumerate}
\def\labelenumi{\arabic{enumi})}
\setcounter{enumi}{2}
\tightlist
\item
  Let B be a ring, a a two-sided ideal of B, b a left ideal of B containing a. Denote by A the ring B/a, by \(p: \mathbf{B} \to \mathbf{A}\) the quotient homomorphism and by \(b'\) the ideal \(p(\mathbf{b}) \subset \mathbf{A}\) . a) Consider the sequence of canonical injections:
\end{enumerate}

\[
\dots \rightarrow a ^ {n} b \rightarrow a ^ {n} \rightarrow a ^ {n - 1} b \rightarrow a ^ {n - 1} \rightarrow \dots \rightarrow a \rightarrow b \rightarrow B.
\]

Show that the sequence obtained by tensor product with B/a

\[
\dots \rightarrow \mathfrak {a} ^ {n} \mathfrak {b} / \mathfrak {a} ^ {n + 1} \mathfrak {b} \rightarrow \mathfrak {a} ^ {n} / \mathfrak {a} ^ {n + 1} \rightarrow \dots \rightarrow \mathfrak {a} / \mathfrak {a} ^ {2} \rightarrow \mathfrak {b} / \mathfrak {a b} \rightarrow \mathbf {A}
\]

defines a left resolution of the A-module A/b'.

\begin{enumerate}
\def\labelenumi{\alph{enumi})}
\setcounter{enumi}{1}
\item
  Suppose that the ideals a and b are free left A-modules. Show that the preceding resolution is a free resolution of the A-module A/b'.
\item
  Let G be a group, k a commutative ring, \(A = \dot{k}^{(G)}\) the algebra of G over k, \(\varepsilon : A \to k\) the homomorphism of k-algebras such that \(\varepsilon(e_g) = 1\) for all \(g \in G\) . One endows k with the structure of A-module (on the left) associated with the homomorphism \(\varepsilon\) . Let \((g_i)_{i \in I}\) be a generating family of G, F the free group constructed on I, \(B = k^{(F)}\) , \(\pi : F \to G\) the homomorphism such that \(\pi(i) = g_i\) for \(i \in I\) , \(p : B \to A\) the homomorphism of k-algebras deduced from \(\pi\) . Show that the ideals a = Ker (p) and b = Ker ( \(\varepsilon \circ p\) ) are free left B-modules; deduce from this a free resolution of the \(k^{(G)}\) -module k.
\end{enumerate}

(According to I, p. 147, exercise 20, the group Ker (π) admits a basic family (rα)α∈J; show that the elements (eᵣα - 1), for α ∈ J (resp. (eᵢ - 1), for i ∈ I) form a basis of the left B-module a (resp. b).)

\begin{enumerate}
\def\labelenumi{\alph{enumi})}
\setcounter{enumi}{3}
\tightlist
\item
  With the notations of c), we set \(\mathbf{R} = \operatorname{Ker}(\pi)\) ; let \(\varphi: \mathbf{R} \to \mathbf{R}/(\mathbf{R}, \mathbf{R})\) be the quotient map. We endow the \(k\) -module \(k \otimes_{\mathbf{Z}} \mathbf{R}/(\mathbf{R}, \mathbf{R})\) of the A-module structure such that
\end{enumerate}

\[
e _ {\pi (f)} \left(1 \otimes \varphi (r)\right) = 1 \otimes \varphi (f r f ^ {- 1})
\]

for all \(r \in \mathbb{R}\) , \(f \in \mathbb{F}\) . Show that there exists an exact sequence of A-modules:

\[
0 \rightarrow k \otimes_ {\mathbf {Z}} \mathrm{R} / (\mathrm{R}, \mathrm{R}) \rightarrow \mathrm{A} ^ {(I)} \rightarrow \mathrm{A} \rightarrow k \rightarrow 0.
\]

\begin{enumerate}
\def\labelenumi{\arabic{enumi})}
\setcounter{enumi}{3}
\tightlist
\item
  Assume the ring A is commutative.
\end{enumerate}

Let G be a group, and K a simplicial scheme (X, p. 177, exercise 19). An action of G on K is a homomorphism from G into the group of bijective simplicial maps from K to itself. a) An action of G on K defines an action of G on the A-module S(K, A) (loc. cit.); show that this action makes S(K, A) into a complex of A \(^{(G)}\) -modules.

\begin{enumerate}
\def\labelenumi{\alph{enumi})}
\setcounter{enumi}{1}
\item
  Assume that G acts freely on the underlying set of K. Show that the A \(^{(G)}\) -modules S \(_{n}\) (K, A) are free.
\item
  Assume that the simplicial scheme K is conical. Show that the \(\mathbf{A}^{(\mathbf{G})}\) -complex \(\mathbf{S}(\mathbf{K},\mathbf{A})\) defines a left resolution of A, equipped with the structure of \(\mathbf{A}^{(\mathbf{G})}\) -module for which \(e_g a = a\) for all \(g \in \mathbf{G}\) , \(a \in \mathbf{A}\) .
\item
  Let \(|\mathbf{G}|\) the simplicial scheme \((\mathbf{G},\mathcal{F})\) , where \(\mathcal{F}\) is the set of finite subsets of G. The group G acts by left translations on the simplicial scheme \(|\mathbf{G}|\) ; show that the \(\mathbf{A}^{(\mathbf{G})}\) -complex S(\textbar G\textbar, A) is a free resolution of the \(\mathbf{A}^{(\mathbf{G})}\) -module A.
\item
  Let \(\mathbf{B}(\mathbf{A}^{(G)},\mathbf{A})\) the standard resolution of the \(\mathbf{A}^{(G)}\) -module A; for \(n\geqslant 0\) , the \(\mathbf{A}^{(G)}\) -module \(\mathbf{B}_n(\mathbf{A}^{(G)},\mathbf{A})\) is identified with \((\mathbf{A}^{(G)})^{\otimes (n + 1)}\) . Show that the A-linear homomorphism \(u:\mathrm{S}(|\mathrm{G}|,\mathrm{A})\to \mathrm{B}(\mathrm{A}^{(G)},\mathrm{A})\) , defined by
\end{enumerate}

\[
u \left(g _ {0}, \dots , g _ {n}\right) = g _ {0} \otimes g _ {0} ^ {- 1} g _ {1} \otimes \dots \otimes g _ {n - 1} ^ {- 1} g _ {n} \quad \text { for } g _ {0}, \dots , g _ {n} \text { in } G,
\]

is an isomorphism of \(A^{(G)}\) -complexes.

\begin{enumerate}
\def\labelenumi{\arabic{enumi})}
\setcounter{enumi}{4}
\tightlist
\item
  \begin{enumerate}
  \def\labelenumii{\alph{enumii})}
  \tightlist
  \item
    Let \(0 \to K \to P \to M \to 0\) , \(0 \to K' \to P' \to M \to 0\) two exact sequences of A-modules, with \(P\) and \(P'\) projective. Show that there exists a commutative diagram:
  \end{enumerate}
\end{enumerate}

\[
\begin{array}{c c c}0 \rightarrow \mathrm{K} \oplus \mathrm{P} ^ {\prime} \rightarrow \mathrm{P} \oplus \mathrm{P} ^ {\prime} \rightarrow \mathrm{M} \rightarrow 0\\\downarrow^ {v}&\downarrow^ {u}&\downarrow^ {l _ {M}}\\0 \rightarrow \mathrm{P} \oplus \mathrm{K} ^ {\prime} \rightarrow \mathrm{P} \oplus \mathrm{P} ^ {\prime} \rightarrow \mathrm{M} \rightarrow 0\end{array}
\]

where the horizontal lines are exact, and where \(u\) and \(v\) are isomorphisms (cf. II, p. 183, exercise 4). b) Let \(n\) be an integer \(\geqslant 0\) , \((\mathbf{P}, p)\) and \((\mathbf{P}', p')\) be two left resolutions of the A-module M, such that \(\mathbf{P}_j = \mathbf{P}_j' = 0\) for \(j > n\) and that the A-modules \(\mathbf{P}_i\) and \(\mathbf{P}_i'\) be projective for \(0 \leqslant i \leqslant n - 1\) . Show that the A-modules \(\bigoplus_{i \geqslant 0} (\mathbf{P}_{2i} \oplus \mathbf{P}_{2i+1}')\) and \(\bigoplus_{i \geqslant 0} (\mathbf{P}_{2j}^{\prime} \oplus \mathbf{P}_{2j+1})\) are isomorphic.

\begin{enumerate}
\def\labelenumi{\alph{enumi})}
\setcounter{enumi}{2}
\tightlist
\item
  Under the assumptions of \(b\) ), show that there exist two projective complexes \(\mathbf{R}\) , \(\mathbf{R}'\) of zero homology, such that \(\mathbf{R}_i = \mathbf{R}_i' = 0\) for \(i \notin [0, n]\) , and an isomorphism of resolutions \(u: \mathbf{P} \oplus \mathbf{R} \to \mathbf{P}' \oplus \mathbf{R}'\) .
\end{enumerate}

¶6). Let M be a left A-module. One calls a presentation of length n, or an n-presentation of M, an exact sequence

\[
\mathrm{L} _ {n} \rightarrow \mathrm{L} _ {n - 1} \rightarrow \dots \rightarrow \mathrm{L} _ {1} \rightarrow \mathrm{L} _ {0} \rightarrow \mathrm{M} \rightarrow 0
\]

where \(L_{i}\) is a free left A-module ( \(0 \leqslant i \leqslant n\) ). The presentation is said to be finite if all the \(L_{i}\) are finitely generated free modules.

If M is a finitely generated left A-module, one denotes by \(\lambda(M)\) the least upper bound (finite or equal to \(+\infty\) ) of the integers \(n \geqslant 0\) such that M admits a finite \(n\) -presentation. If M is not finitely generated, one sets \(\lambda(M) = -1\) .

\begin{enumerate}
\def\labelenumi{\alph{enumi})}
\item
  Let \(0 \to \mathbf{P} \to \mathbf{N} \to \mathbf{M} \to 0\) an exact sequence of left A-modules. One then has \(\lambda(\mathbf{N}) \geqslant \inf (\lambda(\mathbf{P}), \lambda(\mathbf{M}))\) .
\item
  Let \(L_{n} \xrightarrow{d_{n}} L_{n-1} \to \cdots \to L_{0} \to M \to 0\) a finite n-presentation of M; if \(\lambda(M) > n\) show that Ker \((d_{n})\) is a finitely generated A-module (use exercise 5, b)).
\item
  Show that, under the hypotheses of a), one has
\end{enumerate}

\[
\lambda (\mathbf {M}) \geqslant \inf \left(\lambda (\mathbf {N}), \lambda (\mathbf {P}) + 1\right).
\]

\((\mathrm{If} n \leqslant \inf (\lambda(\mathbf{N}), \lambda(\mathbf{P}) + 1)\) , show by induction on \(n\) that \(\lambda(\mathbf{M}) \geqslant n\) reasoning as in \(a\) and using \(b)\) .)

\begin{enumerate}
\def\labelenumi{\alph{enumi})}
\setcounter{enumi}{3}
\tightlist
\item
  Show that, under the hypotheses of a), one has
\end{enumerate}

\[
\lambda (\mathbf {P}) \geqslant \inf \left(\lambda (\mathbf {N}), \lambda (\mathbf {M}) - 1\right).
\]

(Same method as in c).) Deduce that, if \(\lambda(\mathbf{N}) = +\infty\) , then \(\lambda(\mathbf{M}) = \lambda(\mathbf{P}) + 1\) .

\begin{enumerate}
\def\labelenumi{\alph{enumi})}
\setcounter{enumi}{4}
\tightlist
\item
  Deduce from a), c) and d) that, if \(\mathbf{N} = \mathbf{M} \oplus \mathbf{P}\) , we have
\end{enumerate}

\[
\lambda (\mathbf {N}) = \inf \left(\lambda (\mathbf {M}), \lambda (\mathbf {P})\right).
\]

In particular, for N to admit a finite presentation it is necessary and sufficient that the same be true for M and P. f) Let \(N_{1}\) , \(N_{2}\) be two submodules of an A-module M. Suppose that \(N_{1}\) and \(N_{2}\) admit a finite presentation. For \(N_{1} + N_{2}\) to admit a finite presentation, it is necessary and sufficient that \(N_{1} \cap N_{2}\) be finitely generated.

\begin{enumerate}
\def\labelenumi{\arabic{enumi})}
\setcounter{enumi}{6}
\tightlist
\item
  \begin{enumerate}
  \def\labelenumii{\alph{enumii})}
  \tightlist
  \item
    With the notations of exercise 6, show that, if M is a projective module, one has
  \end{enumerate}
\end{enumerate}

\[
\lambda (\mathbf {M}) = - 1 \quad \text { or } \quad \lambda (\mathbf {M}) = + \infty .
\]

If A is a left Noetherian ring, then, for every A-module M, one has \(\lambda(M) = -1\) or \(\lambda(M) = +\infty\) . b) If a is a left ideal of a ring A, which is not finitely generated, \(A_s/a\) is a monogenic A-module which does not admit a finite presentation, in other words \(\lambda(A_s/a) = 0\) (cf. X, p. 11, prop. 6).

\begin{enumerate}
\def\labelenumi{\alph{enumi})}
\setcounter{enumi}{2}
\item
  Give an example of a monogenic left ideal a of a ring A such that \(A_{s}/a\) (which is of finite presentation) admits a dual which is not a finitely generated right A-module.
\item
  Let K be a commutative field, E the vector space \(\mathbf{K}^{(\mathbf{N})}\) , \((e_n)\) the canonical basis of E, T the tensor algebra of E, of which a basis is therefore formed by the finite products \(e_{i_1} e_{i_2} \ldots e_{i_k}\) ( \(k \geqslant 0\) , \(i_j \in \mathbf{N}\) for all \(j\) ). For a given integer \(n\) let b be the two-sided ideal of T generated by the products \(e_1 e_0, e_2 e_1, \ldots, e_n e_{n-1}\) and \(e_{n+k} e_n\) for all \(k \geqslant 1\) ; let A be the quotient ring T/b, and for every integer \(m\) , let \(a_{m}\) the canonical image of \(e_{m}\) in A. Show that, if \(\mathbf{M} = \mathbf{A}_{s}/\mathbf{A}a_{0}\) , we have \(\lambda(\mathbf{M}) = n\) (observe that, for \(m \leqslant n - 1\) the left annihilator of \(a_{m}\) is \(\mathbf{A}a_{m+1}\) and use exercise 6, b)).
\end{enumerate}

\begin{enumerate}
\def\labelenumi{\arabic{enumi})}
\setcounter{enumi}{7}
\tightlist
\item
  Let C be a commutative ring, E, F two C-modules. Show that one has \(\lambda(E \otimes_{C} F) \geqslant \inf(\lambda(E), \lambda(F))\) .
\item
  Let \(\rho: A \to B\) a ring homomorphism, M an A-module.
\end{enumerate}

\begin{enumerate}
\def\labelenumi{\alph{enumi})}
\item
  If B is a flat right A-module and if M admits a finite n-presentation, the B-module B \(\otimes_{\mathbf{A}}\) M admits a finite n-presentation.
\item
  If B is a faithfully flat right A-module (X, p. 169, exercise 9), and if the B-module B ⊗\_A M admits a finite n-presentation, then M admits a finite n-presentation (use exercise 6, b) and exercise 12, p. 169).
\end{enumerate}

\begin{enumerate}
\def\labelenumi{\arabic{enumi})}
\setcounter{enumi}{9}
\tightlist
\item
  Let E be an A-module. One says that E is pseudo-coherent if every finitely generated submodule of E is of finite presentation; every submodule of a pseudo-coherent module is pseudo-coherent. One says that E is coherent if it is pseudo-coherent and finitely generated (hence of finite presentation).
\end{enumerate}

\begin{enumerate}
\def\labelenumi{\alph{enumi})}
\item
  Let \(0 \to E' \to E \to E'' \to 0\) an exact sequence of A-modules. Show that, if E is pseudo-coherent (resp. coherent) and \(E'\) finitely generated, \(E''\) is pseudo-coherent (resp. coherent). Show that, if \(E'\) and \(E''\) are pseudo-coherent (resp. coherent), then so is E. Show that, if E and \(E''\) are coherent, so is \(E'\) (use X, p. 11, prop. 6).
\item
  Let E be a coherent A-module, E' a pseudo-coherent (resp. coherent) A-module. Show that, for every homomorphism \(u: \mathrm{E} \to \mathrm{E}'\) , \(\operatorname{Im}(u)\) and \(\operatorname{Ker}(u)\) are coherent and Coker ( \(u\) ) pseudo-coherent (resp. coherent) (use \(a\) )).
\item
  Show that every direct sum (resp. every finite direct sum) of pseudo-coherent (resp. coherent) modules is a pseudo-coherent (resp. coherent) module.
\item
  If E is a pseudo-coherent module and if M, N are coherent submodules of E, show that \(M + N\) and \(M \cap N\) are coherent (use a) and c)).
\item
  Suppose A is commutative. Show that, if E is a coherent A-module and F a coherent (resp. pseudo-coherent) A-module, Hom\_A(E, F) is a coherent (resp. pseudo-coherent) A-module. (Reduce to the case where F is coherent; consider a finite presentation of E, then use b).)
\end{enumerate}

¶11) a) Let A be a ring. Show that the following properties are equivalent:

\(\alpha)\) The A-module \(\mathbf{A}_s\) is coherent (exercise 10).

β) Every A-module of finite presentation is coherent.

γ) For every finitely presented right A-module M, the dual M* is a finitely generated A-module.

δ) For every set I, the right A-module \(A_{d}^{I}\) is flat.

ε) Any product of flat right A-modules is flat.

(To prove that \(\alpha\) ) implies \(\beta\) ), use Exercise 10, b). To see that \(\delta\) ) implies \(\gamma\) ), use X, p. 14, th. 1; to show that \(\alpha\) ) implies \(\varepsilon\) ), use X, p. 170, exercise 18.)

A ring satisfying these properties is said to be left coherent (or simply coherent), and the notion of a right coherent ring is defined similarly.

\begin{enumerate}
\def\labelenumi{\alph{enumi})}
\setcounter{enumi}{1}
\item
  Show that a left Noetherian ring is coherent. Give an example of a right Artinian ring that is not coherent (cf. VIII, § 1, exercise 4).
\item
  Show that an absolutely flat ring (X, p. 169, exercise 16) is coherent.
\item
  Show that, if A is a coherent ring and M is an A-module, we have \(\lambda(M) = -1\) or \(\lambda(M) = 0\) or \(\lambda(M) = +\infty\) (exercise 6). In particular, every finitely presented A-module admits a resolution by finitely generated free modules.
\item
  Let \((\mathbf{A}_{\alpha}, \varphi_{\beta \alpha})\) an inductive system of rings whose index set is filtered, and let \(\mathbf{A} = \varinjlim \mathbf{A}_{\alpha}\) . We assume that for \(\alpha \leqslant \beta\) , \(\mathbf{A}_{\beta}\) is a flat right \(\mathbf{A}_{\alpha}\) -module. Show that if the \(\mathbf{A}_{\alpha}\) are coherent, so is \(\mathbf{A}\) (Observe that \(\mathbf{A}\) is a flat right \(\mathbf{A}_{\alpha}\) -module for all \(\alpha\) and that for every finitely generated left ideal \(\alpha \subset \mathbf{A}\) there exists an index \(\alpha\) and an ideal \(\alpha_{\alpha}\) of \(\mathbf{A}\) such that \(\alpha\) is isomorphic to \(\mathbf{A} \otimes_{\mathbf{A}} a_{\alpha'}\) .)
\item
  Deduce from e) that every polynomial ring (for any finite or infinite set of indeterminates) over a commutative noetherian ring is coherent. Deduce from this that a quotient ring of a coherent ring is not necessarily coherent.
\item
  For A to be left coherent, it is necessary and sufficient that the left annihilator of every element of A be finitely generated, and that the intersection of two finitely generated left ideals be finitely generated (use exercise 6, f)).
\end{enumerate}

\begin{enumerate}
\def\labelenumi{\arabic{enumi})}
\setcounter{enumi}{11}
\item
  Let c be an infinite cardinal. Suppose that every left ideal of the ring A admits a generating set of cardinality ≤ c. Show that every A-module generated by a set of elements of cardinality ≤ c admits a left resolution by free A-modules of rank ≤ c (cf. VIII, § 1, exercise 13).
\item
  Let M be an A-module. A minimal injective resolution of M is an injective resolution \((e, 1)\) of M such that \(I^n\) is an injective envelope of \(Z^n(I)\) for all \(n \geqslant 0\) .
\end{enumerate}

\begin{enumerate}
\def\labelenumi{\alph{enumi})}
\tightlist
\item
  Every A-module admits minimal injective resolutions; any two such resolutions are isomorphic.
\item
  Let I, I' be two injective resolutions of M, with I minimal; let \(f: I \to I'\) and \(g: I' \to I\) two morphisms of resolutions. Then f is injective, g is surjective, and it is the direct sum of the subcomplexes Im(f) (isomorphic to I) and Ker(g) (with zero homology).
\end{enumerate}

\begin{enumerate}
\def\labelenumi{\arabic{enumi})}
\setcounter{enumi}{13}
\tightlist
\item
  Assume that the ring A is local Noetherian; denote its maximal ideal by m. Let L be a complex of finitely generated A-modules, zero in degrees \textless{} 0, M a finitely generated A-module, and q : L → M a morphism. We make the following hypotheses:
\end{enumerate}

\begin{enumerate}
\def\labelenumi{(\roman{enumi})}
\item
  The homomorphism \(1 \otimes q : (\mathbf{A}/\mathfrak{m}) \otimes_{\mathbf{A}} \mathbf{L}_{0} \to (\mathbf{A}/\mathfrak{m}) \otimes_{\mathbf{A}} \mathbf{M}\) is injective.
\item
  One has \(d_{i}(\mathbf{L}_{i})\subset \mathfrak{m}\mathbf{L}_{i - 1}\) for \(i\geqslant 1\) , and the map
\end{enumerate}

\[
\bar {d} _ {i}: \mathrm{L} _ {i} / \mathrm{mL} _ {i} \rightarrow \mathrm{mL} _ {i - 1} / \mathrm{m} ^ {2} \mathrm{L} _ {i - 1} \quad \text { is injective }.
\]

Let (P, p) be a projective resolution of M, \(u : L \to P\) a morphism such that \(p \circ u = q\) Show that u is injective and that \(u(L)\) is a direct summand of P as a graded A-module (reduce to the case where the resolution P is minimal, and use VIII, § 8, no. 3, cor. 3).

15) Let A be a Noetherian local ring, m its maximal ideal, k = A/m. Let P be a minimal projective resolution of the A-module k; for \(n \geqslant 0\) , we denote by \(b_{n}\) the rank of the free A-module \(P_{n}\) .

\begin{enumerate}
\def\labelenumi{\alph{enumi})}
\item
  Show that we have \(b_0 = 1\) and that \(b_1 = \dim_k (m/m^2)\) is the minimal number of generators of \(m\) .
\item
  If A is commutative, we have \(b_{2} \geqslant \frac{1}{2} b_{1}(b_{1} - 1)\) (consider the exact sequence \(\mathbf{P}_{2} \to m\mathbf{P}_{1} \to m^{2} \to 0\) ). c) Assume the ring A is Artinian (not necessarily commutative). Prove the inequality \(b_{2} \geqslant \frac{1}{4} b_{1}^{2}\) (let \(\mathbf{M}_{l}\) the submodule of \(\mathbf{P}_{2}\) of elements \(x\) such that \(d_{2}(x) \in m^{l}\mathbf{P}_{1}\) ; show that we have
\end{enumerate}

\[
\mathbf {P} _ {2} \subset \mathbf {M} _ {l + 1}
\]

and that there exists an exact sequence \(0 \to \mathbf{M}_l / \mathbf{M}_{l+1} \to \mathfrak{m}^l \mathbf{P}_1 / \mathfrak{m}^{l+1} \mathbf{P}_1 \to \mathfrak{m}^{l+1} / \mathfrak{m}^{l+2} \to 0\) Deduce from this that the polynomial \(p(t) = \sum_{l} (\dim_k (\mathfrak{m}^l / \mathfrak{m}^{l+1})) t^l\) satisfies \((b_2 t^2 - b_1 t + 1)p(t) \geqslant 0\) for \(0 \leqslant t \leqslant 1\) . Conclude

by examining the cases separately \(b_{1} = 1, b_{1} \geqslant 2\) .

\begin{enumerate}
\def\labelenumi{\arabic{enumi})}
\setcounter{enumi}{15}
\tightlist
\item
  \begin{enumerate}
  \def\labelenumii{\alph{enumii})}
  \tightlist
  \item
    Let \(0 \to \mathbf{M}' \to \mathbf{M} \to \mathbf{M}'' \to 0\) an exact sequence of A-modules, \((\mathbf{P}', p')\) (resp. \((\mathbf{P}'', p'')\) ) a projective resolution of \(\mathbf{M}'\) (resp. \(\mathbf{M}''\) ). Show that there exists a projective resolution \((\mathbf{P}, p)\) of \(\mathbf{M}\) and a commutative diagram:
  \end{enumerate}
\end{enumerate}

\[
\begin{array}{c} 0 \to \mathrm{P} ^ {\prime} \to \mathrm{P} \to \mathrm{P} ^ {\prime \prime} \to 0 \\ \downarrow^ {p ^ {\prime}} \quad \downarrow^ {p} \quad \downarrow^ {p ^ {\prime \prime}} \\ 0 \to \mathrm{M} ^ {\prime} \to \mathrm{M} \to \mathrm{M} ^ {\prime \prime} \to 0 \end{array}
\]

where the horizontal lines are exact.

\begin{enumerate}
\def\labelenumi{\alph{enumi})}
\setcounter{enumi}{1}
\tightlist
\item
  Let \(0 \to \mathbf{N}' \to \mathbf{N} \to \mathbf{N}'' \to 0\) a second exact sequence of A-modules, \((\mathbf{Q}, q)\) (resp. \((\mathbf{Q}', q')\) , resp. \((\mathbf{Q}'', q'')\) ) a projective resolution of \(\mathbf{N}\) (resp. \(\mathbf{N}'\) , resp. \(\mathbf{N}''\) ), so that there exists a commutative diagram with exact rows:
\end{enumerate}

\[
\begin{array}{c}0 \rightarrow Q ^ {\prime} \rightarrow Q \rightarrow Q ^ {\prime \prime} \rightarrow 0\\\downarrow^ {q ^ {\prime}} \quad \downarrow^ {q} \quad \downarrow^ {q ^ {\prime \prime}}\\0 \rightarrow N ^ {\prime} \rightarrow N \rightarrow N ^ {\prime \prime} \rightarrow 0\end{array}
\]

Let, on the other hand:

\[
\begin{array}{c} 0 \to \mathrm{M} ^ {\prime} \to \mathrm{M} \to \mathrm{M} ^ {\prime \prime} \to 0 \\ \downarrow^ {u ^ {\prime}} \quad \downarrow^ {u} \quad \downarrow^ {u ^ {\prime \prime}} \\ 0 \to \mathrm{N} ^ {\prime} \to \mathrm{N} \to \mathrm{N} ^ {\prime \prime} \to 0. \end{array}
\]

a commutative diagram of exact sequences, and let \(\tilde{u}':\mathbf{P}'\to\mathbf{Q}'\) and \(\tilde{u}'':\mathbf{P}''\to\mathbf{Q}''\) morphisms of complexes such that \(u'\circ p'=q'\circ\tilde{u}'\) and \(u''\circ p''=q''\circ\tilde{u}'\) . Show that there exists a morphism \(\tilde{u}:\mathbf{P}\to\mathbf{Q}\) such that \(u\circ p=q\circ\tilde{u}\) and that the diagram:

\[
0 \rightarrow \mathrm{P} ^ {\prime} \rightarrow \mathrm{P} \rightarrow \mathrm{P} ^ {\prime \prime} \rightarrow 0\tag{*}
\]

\[
0 \rightarrow Q ^ {\prime} \rightarrow Q \rightarrow Q ^ {\prime \prime} \rightarrow 0
\]

is commutative.

\begin{enumerate}
\def\labelenumi{\alph{enumi})}
\setcounter{enumi}{2}
\tightlist
\item
  Let \(\tilde{u}_{1}^{\prime}: \mathrm{P}^{\prime} \to \mathrm{Q}^{\prime}, \tilde{u}_{1}: \mathrm{P} \to \mathrm{Q}\) and \(\tilde{u}_{1}^{\prime\prime}: \mathrm{P}^{\prime\prime} \to \mathrm{Q}^{\prime\prime}\) three morphisms of complexes such that
\end{enumerate}

\[
u ^ {\prime} \circ p ^ {\prime} = q ^ {\prime} \circ \tilde {u} _ {1} ^ {\prime}, u \circ p = q \circ \tilde {u} _ {1}, u ^ {\prime \prime} \circ p ^ {\prime \prime} = q ^ {\prime \prime} \circ \tilde {u} _ {1} ^ {\prime \prime}
\]

making the diagram (*) commutative; let \(s'\) (resp. \(s''\) ) a homotopy connecting \(\tilde{u}'\) to \(\tilde{u}_{1}'\) (resp. \(\tilde{u}''\) to \(\tilde{u}_{1}'\) ). Show that there exists a homotopy \(s\) connecting \(\tilde{u}\) to \(\tilde{u}_{1}\) , such that the diagram:

\[
\begin{array}{c}0 \rightarrow \mathrm{P} ^ {\prime} \rightarrow \mathrm{P} \rightarrow \mathrm{P} ^ {\prime \prime} \rightarrow 0\\\downarrow s ^ {\prime} \quad \downarrow s \quad \downarrow s ^ {\prime \prime}\\0 \rightarrow \mathrm{Q} ^ {\prime} \rightarrow \mathrm{Q} \rightarrow \mathrm{Q} ^ {\prime \prime} \rightarrow 0\end{array}
\]

is commutative.

\begin{enumerate}
\def\labelenumi{\alph{enumi})}
\setcounter{enumi}{3}
\tightlist
\item
  State and prove the results corresponding to a), b), c) for injective resolutions. ¶17) In this exercise, if \((\mathbf{K}, d', d'')\) is a bicomplex (X, p. 174, exercise 11), we will denote \(K_{p,*}\) the complex \((\bigoplus_{n} K_{p,n}, d'')\) Every complex will be considered as a bicomplex, zero in degree \((p, q)\) for \(q \neq 0\) .
\end{enumerate}

Let C be a complex of A-modules. A projective resolution of C is a pair \((P, p)\) , where P is a bicomplex and \(p : P \to C\) a morphism of bicomplexes, such that \(P_{p,q} = 0\) for q \textless{} 0 and such that the complex \(P_{p,*}\) (resp. \(Z'_{p,*}(\hat{P})\) , \(B'_{p,*}(\mathbf{P})\) , \(H'_{p,*}(\mathbf{P})\) ) define a projective resolution of \(C_p\) (resp. \(Z_p(\mathbf{C})\) , \(B_p(\mathbf{C})\) , \(H_p(\mathbf{C})\) ) for every \(p \in Z\) . For \((P, p)\) be a projective resolution of C, it suffices that the complexes \(B'_{p,*}(\mathbf{P})\) and \(H'_{p,*}(\mathbf{P})\) define, for every p, projective resolutions of \(B_p(\mathbf{C})\) and \(H_p(\mathbf{C})\) .

\begin{enumerate}
\def\labelenumi{\alph{enumi})}
\tightlist
\item
  Show that every A-complex C admits a projective resolution (choose, for every p, projective resolutions \(B_{p,*}\) and \(H_{p,*}\) of \(B_{p}(C)\) and \(H_{p}(C)\) ; using Exercise 16, construct projective resolutions \(Z_{p,*}\) of \(Z_{p}(C)\) and \(P_{p,*}\) of \(C_{p}\) as well as exact sequences:
\end{enumerate}

\[
0 \rightarrow \mathrm{B} _ {p, *} \rightarrow \mathrm{Z} _ {p, *} \rightarrow \mathrm{H} _ {p, *} \rightarrow 0; \quad 0 \rightarrow \mathrm{Z} _ {p, *} \rightarrow \mathrm{P} _ {p, *} \rightarrow \mathrm{B} _ {p + 1, *} \rightarrow 0).
\]

\begin{enumerate}
\def\labelenumi{\alph{enumi})}
\setcounter{enumi}{1}
\item
  Let \(\mathbf{C}'\) another complex, \((\mathbf{P}', p')\) a projective resolution of \(\mathbf{C}'\) , \(u: \mathbf{C} \to \mathbf{C}'\) a morphism of complexes. Show that there exists a morphism of bicomplexes \(\tilde{u}: \mathbf{P} \to \mathbf{P}'\) such that \(u \circ p = p' \circ \tilde{u}\) (for every \(p \in \mathbf{Z}\) , choose morphisms \(\mathbf{B}_{p'}', (\mathbf{P}) \to \mathbf{B}_{p'}', (\mathbf{P}')\) and \(\mathrm{H}_{p'}', (\mathbf{P}) \to \mathrm{H}_{p'}', (\mathbf{P}')\) ; construct \(\tilde{u}\) as before, using Exercise 16, b)).
\item
  Let \(v: \mathbf{C} \to \mathbf{C}'\) a morphism of complexes homotopic to \(u\) , and let \(\tilde{v}: \mathbf{P} \to \mathbf{P}'\) be a morphism of bicomplexes such that \(v \circ p = p' \circ \tilde{v}\) . Show that \(\tilde{u}\) and \(\tilde{v}\) are homotopic (in the sense of X, p. 174, Exercise 11, c); use Exercise 16, c)).
\end{enumerate}

\(d)\) Let E be a subset of \(\mathbf{Z}\) , and C an A-complex such that \(\mathbf{C}_{p}=0\) for \(p\in\mathbf{E}\) . Show that there exists a projective resolution \((\mathbf{P},p)\) of C such that \(\mathbf{P}_{p,q}=0\) for \(p\in\mathbf{E}\) .

\begin{enumerate}
\def\labelenumi{\alph{enumi})}
\setcounter{enumi}{4}
\item
  Suppose there exists an integer \(n\) such that every A-module admits a projective resolution of length \(\leqslant n\) . Show that every A-complex C admits a projective resolution \((P, p)\) such that \(P_{p,q} = 0\) for \(q > n\) .
\item
  Let C be an A-complex. An injective resolution of C is a pair \((e, l)\) , where I is a bicomplex and \(e: C \to 1\) a morphism of bicomplexes, such that \(I^{p,q} = 0\) for q \textless{} 0 and such that the complex \(I^{p,*}\) (resp. \(Z^{p,*}(I)\) , \(B^{p,*}(l)\) , \(H^{p,*}(l)\) ) defines an injective resolution of \(C^{p}\) (resp. \(Z^{p}(C)\) , \(B^{p}(C)\) , \(H^{p}(C)\) ) for every \(p \in Z\) . State and prove the analogues of results a) to e) for injective resolutions.
\end{enumerate}

\begin{enumerate}
\def\labelenumi{\arabic{enumi})}
\setcounter{enumi}{17}
\tightlist
\item
  Assume the ring A is commutative. A differential graded A-algebra is an A-algebra S, graded of type Z, such that \(S_{n} = 0\) for n \textless{} 0, equipped with an antiderivation d of degree (-1), satisfying \(d \circ d = 0\) . We shall also denote by S the complex of A-modules (S, d).
\end{enumerate}

\begin{enumerate}
\def\labelenumi{\alph{enumi})}
\item
  Let C be a complex of A-modules, zero in degrees \textless{} 0. Show that there exists on S = T(C) (resp. S = S(C), resp. S = Λ(C)) a structure of differential graded A-algebra, such that the canonical injection of C into S is a morphism of complexes.
\item
  Show that the multiplication of S induces a structure of graded A-algebra on H(S).
\item
  Let S, T be two differential graded A-algebras. Set D(s ⊗ t) = ds ⊗ t + (-1)\^{}p s ⊗ dt for s ∈ S\_p, t ∈ T. Show that D defines a structure of differential graded A-algebra on the left tensor product S *⊗\_A T (III, p. 49).
\end{enumerate}

\begin{enumerate}
\def\labelenumi{\arabic{enumi})}
\setcounter{enumi}{18}
\tightlist
\item
  Assume the ring A is commutative.
\end{enumerate}

\begin{enumerate}
\def\labelenumi{\alph{enumi})}
\item
  Let \((\mathbb{S}, d_{\mathbb{S}})\) a differential graded A-algebra (exercise 18) , M an A-module, \(u: \mathbb{M} \to Z_{n}(\mathbb{S})\) an A-homomorphism. Show that there exists on \(\mathbb{S} \otimes_{\mathbb{A}} \mathbf{T}(\mathbb{M})\) a structure of differential graded A-algebra such that \(d(s \otimes 1) = d_{\mathbb{S}} s \otimes 1\) for \(s \in \mathbb{S}\) and \(d(1 \otimes m) = u(m) \otimes 1\) for \(m \in \mathbb{M}\) .
\item
  Let B be an A-algebra. Show that there exists a differential graded A-algebra S and a homomorphism of A-algebras \(p: \mathbf{S}_{0} \to \mathbf{B}\) such that \((\mathbf{S}, p)\) be a free resolution of the A-module B. (Construct by induction on \(n\) , using \(a\) ), a differential graded algebra \(\mathbf{S}(n)\) , free in each degree, such that \(\mathrm{H}_{i}(\mathbf{S}(n)) = 0\) for \(0 < i < n\) and \(\mathrm{H}_{0}(\mathbf{S}(n)) = \mathbf{B}\) .)
\item
  If the algebra B is commutative, show that one can find an alternating A-algebra S (III, p. 53) satisfying the conditions of b). If moreover A is noetherian and if B is a quotient of A by an ideal, show that one can find S such that \(S_{0} = A\) and that \(S_{n}\) is a finitely generated A-module for every n.
\end{enumerate}

¶ 20) Assume that the ring A is an algebra over a commutative ring k. Let \(\eta: k \to A\) be the homomorphism defined by \(\eta(\lambda) = \lambda 1_{A}\) for \(\lambda \in k\) , and let \(\overline{A}\) its cokernel. If \(a\) is an element of A, we denote by \(\overline{a}\) its class in \(\overline{A}\) . We set \(\mathrm{N}_{n}(\mathrm{A}) = \mathrm{A} \otimes_{k} \overline{\mathrm{A}}^{\otimes n} \otimes_{k} \mathrm{A}\) for \(n \geqslant 0\) , \(\mathrm{N}_{n}(\mathrm{A}) = 0\) for \(n < 0\) , and

\[
\mathrm{N} (\mathbf {A}) = \bigoplus_ {n \in \mathbf {Z}} \mathrm{N} _ {n} (\mathbf {A}).
\]

\begin{enumerate}
\def\labelenumi{\alph{enumi})}
\tightlist
\item
  Show that there exists an (A, A)-endomorphism d of N(A), graded of degree (-1), such that
\end{enumerate}

\[
\begin{array}{l} d (a \otimes \bar {a} _ {1} \otimes \dots \otimes \bar {a} _ {n} \otimes b) = a a _ {1} \otimes \bar {a} _ {2} \otimes \dots \otimes \bar {a} _ {n} \otimes b + \sum_ {i = 1} ^ {n - 1} (- 1) ^ {i} a \otimes \dots \otimes \overline {{a _ {i}}} \overline {{a _ {i + 1}}} \otimes \dots \otimes \bar {a} _ {n} \otimes b \\ \qquad + (- 1) ^ {n} a \otimes \bar {a} _ {1} \otimes \dots \otimes \bar {a} _ {n - 1} \otimes a _ {n} b \quad \text {for} \quad a, a _ {1},..., a _ {n}, b \in A \end{array}
\]

\begin{enumerate}
\def\labelenumi{\alph{enumi})}
\setcounter{enumi}{1}
\item
  Show that \(d \circ d = 0\) , so that \(\mathbf{N}(\mathbf{A})\) is a complex of \((\mathbf{A}, \mathbf{A})\) -bimodules. If M is a left A-module, we denote by \(\mathbf{N}(\mathbf{A}, \mathbf{M})\) the A-complex \(\mathbf{N}(\mathbf{A}) \otimes_{\mathbf{A}} \mathbf{M}\) .
\item
  Let B(A, M) be the standard resolution of M (X, p. 58), \(p: \mathrm{B}(\mathrm{A}, \mathrm{M}) \to \mathrm{N}(\mathrm{A}, \mathrm{M})\) the A-homomorphism defined by \(p(a_0 \otimes \cdots \otimes a_n \otimes m) = a_0 \otimes \overline{a}_1 \otimes \cdots \otimes \overline{a}_n \otimes m\) for \(a_0, \ldots, a_n \in \mathbf{A}, m \in \mathbf{M}\) . Show that \(p\) is a homotopy equivalence of A-complexes (construct, by induction, a map inverse up to homotopy).
\end{enumerate}

* 21) Let \(k\) be a commutative ring, and let g be a \(k\) Lie algebra, U its enveloping algebra (cf. LIE, I, § 2). Assume that g is a \(k\) -free module.

\begin{enumerate}
\def\labelenumi{\alph{enumi})}
\tightlist
\item
  Show that for every \(n \geqslant 0\) an injective U-homomorphism
\end{enumerate}

\[
j _ {n}: \mathrm{U} \otimes_ {k} \Lambda^ {n} (\mathrm{g}) \rightarrow \mathrm{U} ^ {\otimes (n + 1)} \quad \text { such that } \quad j _ {n} (u \otimes (x _ {1} \wedge \dots \wedge x _ {n})) = \sum_ {\sigma \in \mathfrak {S} _ {n}} \varepsilon (\sigma), u \otimes x _ {\sigma (1)} \otimes \dots \otimes x _ {\sigma (n)}
\]

for \(u\in \mathbf{U},x_1,\dots ,x_n\in \mathfrak{g}\)

We thus define a graded U-homomorphism of degree 0. \(j:U\otimes_{k}\Lambda(g)\to B(U,k)\) , where \(B(U,k)\) is the standard resolution of the left U-module \(k\) (the structure of U-module of \(k\) being deduced from the natural homomorphism \(U\to k\) ). Show that \(U\otimes_{k}\Lambda(g)\) is identified by \(j\) with a subcomplex of \(B(U,k)\) , the differential being given by the formula :

\[
\begin{array}{l} d (u \otimes (x _ {1} \wedge \dots \wedge x _ {n})) = \sum_ {1 \leqslant i \leqslant n} (- 1) ^ {i + 1} u x _ {i} \otimes (x _ {1} \wedge \dots \wedge \hat {x} _ {i} \wedge \dots \wedge x _ {n}) \\ + \sum_ {1 \leqslant i <   j \leqslant n} (- 1) ^ {i + j} u \otimes ([ x _ {i}, x _ {j} ] \wedge x _ {1} \wedge \dots \wedge \hat {x} _ {i} \wedge \dots \wedge \hat {x} _ {j} \wedge \dots \wedge x _ {n}) \quad \text { pour } \quad u \in U, x _ {1},..., x _ {n} \in g. \end{array}
\]

(where the sign over a letter means that it is to be omitted). We denote by V(g) the U-complex thus defined. b) Show that V(g) defines a free resolution of the U-module k (let \((\mathrm{U}_{n})_{n\geqslant0}\) be the natural filtration of U, \(F_{r}V(g)\) the sub-k-module of V(g) generated by the elements \(u_{p}\otimes x_{q}\) for \(u_{p}\in U_{p}, x_{q}\in\Lambda^{q}(g), p+q\leqslant r\) . Show that \(F_{r}V(g)\) is a subcomplex of V(g), and that the complex \(\bigoplus_{r\geqslant0}F_{r}V(g)/F_{r-1}V(g)\)

is isomorphic to the complex \(\mathbf{S}(\mathfrak{g})\otimes\symbf{\Lambda}(\mathfrak{g})\) defined in X, p. 151). *
% semantic-fragments: legacy-input-seam
\relax{}
